\section{\textit{Quasinormal Modes} (QNM)}
\label{sec:qnm}

\noindent Sistem fisis yang mengalami gangguan kecil di sekitar keadaan setimbangnya pada
umumnya berosilasi pada sekumpulan frekuensi khas yang disebut mode normal. Contoh paling
sederhana dijumpai pada dawai yang kedua ujungnya terikat: gangguan pada dawai tersebut
terurai menjadi gelombang-gelombang tegak dengan frekuensi diskret yang ditentukan oleh
panjang dawai beserta syarat batas di kedua ujungnya. Dalam sistem konservatif semacam ini,
energi gangguan tidak dapat keluar dari sistem, sehingga tiap mode berosilasi dengan
frekuensi yang bernilai real dan amplitudonya bertahan tanpa teredam. Faktor evolusi waktu
mode normal berbentuk $e^{-i\omega t}$ dengan $\omega$ real, sehingga hanya menghasilkan
osilasi murni tanpa peluruhan.

Situasi berbeda terjadi pada sistem yang bersifat terbuka atau berdisipasi, yaitu sistem
yang dapat kehilangan energi ke lingkungan sekitarnya. Contoh paling sederhana adalah
osilator teredam, yang mengalami gaya gesek sebanding kecepatan sehingga energinya berkurang
seiring waktu \parencite{boas2006}. Contoh lain yang lebih dekat dengan persoalan gelombang
adalah rongga akustik atau optik yang temboknya tidak sepenuhnya memantulkan, sehingga
sebagian gelombang di dalamnya lolos keluar. Pada sistem semacam ini, syarat batas tidak lagi
mengurung energi sepenuhnya, sehingga frekuensi karakteristiknya tidak lagi bernilai real,
melainkan kompleks. Mode osilasi teredam semacam ini secara umum disebut
\textit{quasinormal modes}, yakni mode yang menyerupai mode normal namun dengan frekuensi
kompleks akibat sifat terbuka dari sistem yang ditinjau \parencite{Kokkotas1999}.

Frekuensi kompleks tersebut dapat dituliskan secara sederhana sebagai
\begin{equation}
    \omega = \omega_R - i\omega_I,
    \label{eq:omega-qnm}
\end{equation}
dengan $\omega_R$ bagian real yang menyatakan laju osilasi dan $\omega_I$ bagian imajiner
yang menyatakan laju redaman ($\omega_I > 0$). Spektrum QNM terdiri atas deretan mode yang
diberi label indeks \textit{overtone} $n = 0, 1, 2, \dots$, dengan nilai $n$ yang lebih besar
pada umumnya menghasilkan $\omega_I$ yang lebih besar pula, sehingga mode dengan
\textit{overtone} lebih tinggi meluruh lebih cepat. Mode dengan redaman paling lambat, yaitu
$\omega_I$ terkecil pada $n = 0$, disebut mode fundamental. Sebagai gambaran, hal ini
menyerupai lonceng yang dipukul: nada dasarnya bertahan berdengung paling lama, sedangkan
nada-nada tambahan di atasnya lebih cepat meredup, sehingga bunyi yang masih terdengar
beberapa saat kemudian didominasi oleh nada dasar tersebut.

Lubang hitam yang dikenai gangguan merupakan salah satu wujud sistem terbuka tersebut.
Energi gangguan tidak terkurung, melainkan dapat lolos dengan cara terserap masuk ke horizon
peristiwa maupun teradiasikan menuju tak hingga spasial, sehingga osilasinya teredam seiring
waktu \parencite{Regge1957, Vishveshwara1970}. Frekuensi quasinormal lubang hitam diperoleh
sebagai solusi persamaan diferensial \eqref{eq:master-skalar} dengan nilai frekuensi
berbentuk Persamaan~\eqref{eq:omega-qnm}. Sejalan dengan itu, frekuensi mode kuasinormal
dapat digunakan untuk menganalisis kestabilan lubang hitam. Lubang hitam dikatakan stabil
terhadap suatu jenis perturbasi apabila nilai $\omega_I$ bernilai positif, sehingga setiap
gangguan pasti akan teredam seiring berjalannya waktu. Sebaliknya, apabila $\omega_I$ bernilai
negatif, gangguan justru tumbuh secara eksponensial, sehingga lubang hitam dikatakan tidak
stabil, yaitu gangguan yang amplitudonya membesar tanpa henti tanpa mengalami peredaman
\parencite{Konoplya2011}.

Secara matematis, peluruhan maupun pertumbuhan gangguan tersebut dapat dilihat dari faktor
evolusi waktu. Solusi persamaan master \eqref{eq:master-skalar} hanya memuat informasi bagian
radial, sehingga perlu digabungkan kembali dengan faktor evolusi waktu $e^{-i\omega t}$ yang
telah dipisahkan melalui ansatz \eqref{eq:ansatz}, menghasilkan bentuk gelombang utuh
\begin{equation}
    \phi \sim e^{-i\omega t}e^{\pm i\omega r_*}.
    \label{eq:gelombang-utuh}
\end{equation}
Substitusi $\omega = \omega_R - i\omega_I$ ke dalam faktor waktu memberikan
$e^{-i\omega t} = e^{-i\omega_R t}\, e^{-\omega_I t}$, yang memperlihatkan bahwa
$\omega_I > 0$ menghasilkan peluruhan eksponensial sedangkan $\omega_I < 0$ menghasilkan
pertumbuhan eksponensial, sesuai dengan interpretasi kestabilan yang telah diuraikan di atas.
Selain faktor waktu, bentuk gelombang utuh \eqref{eq:gelombang-utuh} juga memperlihatkan arah
rambat gelombang di kedua ujung domain melalui tanda pada eksponen $e^{\pm i\omega r_*}$. Dengan
arah rambat tersebut, kondisi batas fisik yang diberlakukan di horizon peristiwa serta di tak
hingga spasial dapat dinyatakan, sehingga frekuensi eigen dapat diperoleh secara lengkap.

Potensial efektif sangat berpengaruh dari kondisi batas ruang-waktu yang ditinjau. Pada ruang-waktu
Schwarzschild yang datar secara asimtotik, potensial efektif meluruh menuju nol pada kedua
ujung domain tersebut. Akibatnya, suku potensial pada persamaan master
\eqref{eq:master-skalar} dapat dianggap nol dan persamaan tereduksi menjadi
\begin{equation}
    \frac{d^2\Psi}{dr_*^2} + \omega^2\Psi = 0
\end{equation}
yang memiliki solusi umum,
\begin{equation}
    \Psi \sim e^{\pm i\omega r_*}
\end{equation}
dalam koordinat \textit{tortoise} \parencite{Konoplya2011, Horowitz2000}. Solusi ini
merupakan bagian radial dari bentuk gelombang utuh \eqref{eq:gelombang-utuh}, sehingga arah
rambat tiap solusi dapat ditinjau dengan menggabungkannya kembali dengan faktor evolusi
waktu. Digabungkan dengan faktor tersebut, solusi $e^{-i\omega r_*}$ merambat ke arah $r_*$
mengecil dan diinterpretasikan sebagai gelombang masuk murni di horizon peristiwa, sedangkan
solusi $e^{+i\omega r_*}$ merambat ke arah $r_*$ membesar dan diinterpretasikan sebagai
gelombang keluar murni menuju daerah asimtotik. Dengan demikian, kondisi batas fisik untuk
lubang hitam Schwarzschild memberlakukan gelombang masuk murni di horizon peristiwa dan
gelombang keluar murni di tak hingga spasial,
\begin{align}
    \Psi &\sim e^{-i\omega r_*}, \qquad r_* \to -\infty \quad (\text{horizon peristiwa}), \notag \\
    \Psi &\sim e^{+i\omega r_*}, \qquad r_* \to +\infty \quad (\text{tak hingga spasial}).
    \label{eq:bc-schwarzschild}
\end{align}
Kondisi batas di horizon merepresentasikan bahwa energi jatuh ke dalam lubang hitam tanpa
pantulan, sedangkan kondisi batas di tak hingga merepresentasikan bahwa hanya radiasi yang
menjauhi lubang hitam yang diperbolehkan.

Situasi berbeda terjadi pada latar belakang Schwarzschild Anti-de Sitter. Kehadiran suku
tambahan dalam fungsi metrik mengubah perilaku dan karakteristik geometris, sehingga
potensial efektif tidak meluruh menuju nol di tak hingga, melainkan tumbuh tanpa batas
($V(r) \to \infty$ ketika $r \to \infty$), sehingga tak hingga spasial berperilaku sebagai
dinding pemantul. Berbeda dengan kasus Schwarzschild yang dapat kehilangan energi melalui horizon peristiwa
maupun tak hingga spasial, pada latar belakang ini energi tidak dapat lolos menuju tak
hingga, sehingga satu-satunya cara gangguan meluruh adalah dengan terserap masuk ke horizon
peristiwa. Mekanisme ini sudah cukup untuk membuat frekuensi tetap bernilai kompleks,
sehingga gangguan tetap meluruh secara eksponensial meskipun energi tidak dapat lolos menuju
tak hingga. Kondisi batas gelombang keluar murni tidak dapat diterapkan seperti pada
Schwarzschild karena gelombang divergen saat menyentuh batas tak hingga spasial anti-de
Sitter. Untuk memenuhi kondisi batas fisik, diberlakukan syarat tambahan pada tak hingga
spasial, yaitu gelombang harus lenyap \parencite{Horowitz2000}. Dengan demikian, kondisi
batas untuk lubang hitam Schwarzschild anti-de Sitter memberlakukan gelombang masuk murni
di horizon peristiwa dan kondisi Dirichlet di tak hingga spasial,
\begin{alignat}{3}
    \Psi &\sim e^{-i\omega r_*}, &\quad &r_* \to -\infty \quad &&(\text{horizon peristiwa}), \notag \\
    \Psi &\to 0, &\quad &r_* \to \text{constant} \quad &&(\text{batas anti-de Sitter}).
    \label{eq:bc-sads}
\end{alignat}
Kondisi yang melenyapkan medan pada batas anti-de Sitter ini dikenal sebagai kondisi batas
Dirichlet, yang memilih mode \textit{normalizable} dan konsisten dengan kerangka
korespondensi anti-de Sitter \parencite{MossNorman2002}.